% TOP_PROBLEM: {"id": 3000018, "title": "Complexity of the halting problem for simple digraphs", "category_id": 2, "category": {"id": 2, "name": "combinatorics", "display_name": "Combinatorics", "description": "Counting problems, graph theory, discrete structures.", "slug": "combinatorics", "order_index": 2, "created_at": "2026-07-31T15:26:25.671Z"}, "status": "open"}
% FIRST_POSTED: null
% ENTRY_KIND: statement_only
% Imported from: batch09.tex; UnsolvedMath ID 3000018
\documentclass[11pt]{article}
\usepackage[margin=25mm]{geometry}
\usepackage{fontspec}
\setmainfont{Latin Modern Roman}
\usepackage{amsmath,amssymb,mathtools}
\usepackage{unicode-math}
\setmathfont{Latin Modern Math}
\usepackage{xurl,hyperref}
\hypersetup{colorlinks=true,urlcolor=blue}
\setcounter{secnumdepth}{0}
\setlength{\parindent}{0pt}
\setlength{\parskip}{5pt}
\setlength{\emergencystretch}{3em}
\newcommand{\sref}[2]{\hyperlink{source:#1:#2}{[S#2]}}
\newcommand{\cataloguescope}[1]{\paragraph{Recorded target.}#1}
\begin{document}
% UnsolvedMath ID 3000018; source catalogue code AMR-029-0018
\setcounter{section}{3000017}
\section{Complexity of the halting problem for simple digraphs}\label{3000018}
{\small\sffamily UnsolvedMath ID 3000018 · Combinatorics}
\subsection{Definitions and mathematical statement}

\paragraph{Shared notation.}$\mathbb N=\{1,2,\ldots\}$, and $\mathbb Z,\mathbb Q,\mathbb R,\mathbb C$ denote the integers, rationals, reals and complex numbers. Any other conventions are specified locally.


Let $G=(V,A)$ be a finite loopless directed multigraph, with $a_{uv}$ arcs from $u$ to $v$ and $d^+(v)=\sum_w a_{vw}$. A chip distribution is $x\in\mathbb Z_{\ge0}^{V}$. Firing $v$ is legal when $x(v)\ge d^+(v)$ and changes $x$ by removing $d^+(v)$ chips at $v$ and sending one chip along each outgoing arc. Reachability permits any finite legal firing sequence, including the empty sequence. A distribution is terminating if no infinite legal firing sequence starts from it; equivalently a maximal legal game reaches a distribution with no legal firing. The original legal-firing convention includes a zero-outdegree vertex, whose firing has no effect. Arc multiplicities and chip counts are binary encoded. A simple digraph here has no loops and at most one arc for each ordered pair of distinct vertices; opposite arcs are allowed. The halting decision problem asks whether its given chip distribution is terminating.
\paragraph{Statement recorded in the supplied source.}
Is the chip-firing halting decision problem NP-complete for simple digraphs? \sref{3000018}{2}
\subsection{Short English statement}
Is deciding chip-firing termination NP-complete even when there are no parallel arcs?
\subsection{Sources}
\begin{description}
\item[\hypertarget{source:3000018:1}{S1}] ULAM.AI and the UnsolvedMath Contributors, \emph{UnsolvedMath: A Curated Collection of Open Mathematics Problems}, record 3000018 (AMR-029-0018). CC BY 4.0. \url{https://huggingface.co/datasets/ulamai/UnsolvedMath}.
\item[\hypertarget{source:3000018:2}{S2}] B. Hujter, V. Kiss and L. Tóthmérész, On the complexity of the chip-firing reachability problem (2017), Sections 1 and 2.1--2.4; Section 5, Problems 19 and 22, pp.12--13. \url{https://arxiv.org/pdf/1507.03209}.
\end{description}
\label{end:3000018}
\end{document}
